\documentclass[14pt,a4paper]{article}
\usepackage{geometry}
\geometry{left=3cm, right=1.5cm, top=2cm, bottom=2cm}

\usepackage{fontspec}
\setmainfont{Times New Roman}
\usepackage[russian]{babel}
\usepackage{indentfirst}
\usepackage{enumitem}
\usepackage{booktabs}
\usepackage{hyperref}
\hypersetup{colorlinks=true, linkcolor=black, citecolor=black, urlcolor=black}

\begin{document}

%==================== ТИТУЛЬНЫЙ ЛИСТ ====================
\begin{titlepage}
\centering
\textbf{МИНИСТЕРСТВО СЕЛЬСКОГО ХОЗЯЙСТВА РОССИЙСКОЙ ФЕДЕРАЦИИ}\\[0.8cm]

\textbf{ФЕДЕРАЛЬНОЕ ГОСУДАРСТВЕННОЕ БЮДЖЕТНОЕ ОБРАЗОВАТЕЛЬНОЕ УЧРЕЖДЕНИЕ ВЫСШЕГО ОБРАЗОВАНИЯ}\\[0.6cm]

\textbf{«РОССИЙСКИЙ ГОСУДАРСТВЕННЫЙ АГРАРНЫЙ УНИВЕРСИТЕТ – МСХА имени К.А. ТИМИРЯЗЕВА»}\\[0.4cm]
(ФГБОУ ВО РГАУ - МСХА имени К.А. Тимирязева)\\[1cm]

Институт экономики и управления АПК\\
Кафедра педагогики и психологии профессионального образования\\[1.5cm]

\textbf{КУРСОВАЯ РАБОТА}\\[1cm]

По дисциплине: «Информационные технологии в профессиональном образовании»\\[0.8cm]

\textbf{ТЕМА:}\\[0.3cm]
«Инструмент для составления расписаний, учебных планов и напоминаний »\\[1.5cm]

Направление подготовки: 44.03.04 Профессиональное обучение (по отраслям)\\
Направленность: Экономика и управление, Информационные технологии в образовании\\[1.2cm]

\flushright
Выполнила\\
студентка 2 курса\\
группы ДЭ 21-25\\
Татаринова Александра Геннадьевна\\[0.8cm]

Дата регистрации КР на кафедре:\\
\rule{7cm}{0.4pt}\\[0.5cm]

Допущен(а) к защите:\\
\rule{7cm}{0.4pt}\\[0.8cm]

Руководитель:\\
к.п.н., доцент, Заифирова Л.В.\\
(ученая степень, ученое звание, ФИО)\\
\rule{7cm}{0.4pt}\\
подпись\\[0.8cm]

Члены комиссии:\\
\rule{7cm}{0.4pt}\\
(ученая степень, ученое звание, ФИО)\qquad подпись\\[0.3cm]
\rule{7cm}{0.4pt}\\
(ученая степень, ученое звание, ФИО)\qquad подпись\\[0.3cm]
\rule{7cm}{0.4pt}\\
(ученая степень, ученое звание, ФИО)\qquad подпись\\[0.8cm]

Оценка: \rule{5cm}{0.4pt}\\
Дата защиты: \rule{5cm}{0.4pt}\\[1.5cm]

\centering
Москва, 2025
\end{titlepage}


%==================== ЗАДАНИЕ НА КУРСОВУЮ РАБОТУ ====================
\newpage
\centering
\textbf{МИНИСТЕРСТВО СЕЛЬСКОГО ХОЗЯЙСТВА РОССИЙСКОЙ ФЕДЕРАЦИИ}\\[0.4cm]
Российский государственный аграрный университет – МСХА имени К.А. Тимирязева\\
Институт экономики и управления АПК\\
Кафедра педагогики и психологии профессионального образования\\[1cm]

\textbf{ЗАДАНИЕ НА КУРСОВУЮ РАБОТУ (КР)}\\[1cm]

\noindent Студент: Лю Янь\\
Группа: Калитвина\\[0.4cm]

Тема КР: «Инструмент для составления расписаний, учебных планов и напоминаний как средство повышения учебной самоорганизации студентов профессиональных образовательных организаций»\\[0.6cm]

Исходные данные к работе: научная и учебная педагогическая литература, периодические издания по цифровизации образования, интернет-ресурсы, существующие программные решения для планирования учебной деятельности, результаты анкетирования студентов СПО.\\[0.6cm]

Перечень подлежащих разработке в работе вопросов:
\begin{enumerate}
    \item Теоретическое изучение сущности учебной самоорганизации студентов СПО, особенностей использования цифровых инструментов в образовательном процессе профессиональных образовательных организаций.
    \item Анализ проблем самоорганизации и планирования учебной деятельности у студентов среднего профессионального образования.
    \item Разработка методики применения цифрового инструмента для составления расписаний, учебных планов и напоминаний в условиях ПОО.
    \item Оценка эффективности использования данного инструмента для повышения уровня учебной самоорганизации студентов СПО.
    \item Разработка рекомендаций по адаптации и дальнейшему внедрению инструмента в учебный процесс профессиональных образовательных организаций.
\end{enumerate}

Перечень дополнительного материала: описание функционала инструмента, пример шаблона учебного плана и расписания, бланки анкет для входного и итогового опроса, презентация-сопровождение для занятия со студентами.\\[1cm]

Дата выдачи задания «\_\_\_»\_\_\_\_\_\_\_\_202\_\_ г.\\[0.4cm]
Руководитель (подпись, ФИО) \rule{6cm}{0.4pt}\\[0.4cm]
Задание принял к исполнению (подпись студента) \rule{6cm}{0.4pt}\\
«\_\_\_»\_\_\_\_\_\_\_\_202\_\_ г.

%==================== СОДЕРЖАНИЕ ====================
\newpage
\tableofcontents

%==================== ВВЕДЕНИЕ ====================
\newpage
\section*{Введение}
\addcontentsline{toc}{section}{Введение}

Большинство студентов среднего профессионального образования обладают хорошими способностями к освоению профессиональных дисциплин, но при этом часто испытывают существенные трудности с самостоятельным планированием собственного учебного времени. Неправильное распределение часов на подготовку к занятиям, забывание о сроках сдачи курсовых работ, лабораторных заданий и практик приводят к пропускам, просрочкам, снижению академической успеваемости и потере учебной мотивации. Поэтому формирование навыков учебной самоорганизации нельзя рассматривать как второстепенную задачу образовательного процесса; оно является важнейшим компонентом подготовки конкурентоспособного специалиста среднего звена.

Современный рынок труда предъявляет высокие требования не только к профессиональным знаниям и умениям выпускников СПО, но и к их личностным качествам: умению планировать рабочее время, расставлять приоритеты, своевременно выполнять поставленные задачи. Данные компетенции формируются уже в период обучения в колледже. В настоящее время широкое распространение получают цифровые инструменты, которые могут выступать вспомогательным средством для организации учебной деятельности. К ним относятся сервисы для составления расписаний, формирования индивидуальных учебных планов и системы напоминаний о важных учебных событиях.

Актуальность темы курсовой работы обусловлена потребностью поиска доступных, простых в освоении цифровых средств, направленных на совершенствование навыков самоорганизации обучающихся СПО. Несмотря на большое количество существующих программных решений, во многих профессиональных образовательных организациях отсутствуют четкие методики их целенаправленного применения в воспитательной и учебной работе. Часто студенты используют цифровые сервисы хаотично, без понимания принципов рационального планирования времени. Особенно острая проблема наблюдается среди первокурсников, которые переходят от школьной системы строгого внешнего контроля к самостоятельной организации собственной учебной деятельности в колледже.

Основная цель курсовой работы заключается в изучении педагогических возможностей инструмента для составления расписаний, учебных планов и напоминаний как средства повышения учебной самоорганизации студентов профессиональных образовательных организаций.

Задачи курсовой работы:
\begin{enumerate}
    \item Изучить сущность, структуру и содержание понятия «учебная самоорганизация» студентов СПО.
    \item Проанализировать теоретические основы применения цифровых инструментов планирования в системе среднего профессионального образования.
    \item Выявить психолого-педагогические особенности и типичные трудности самоорганизации учебной деятельности у студентов СПО.
    \item Разработать методику практического занятия, посвященного освоению инструмента для составления расписаний, учебных планов и напоминаний для обучающихся СПО.
\end{enumerate}

Объект курсовой работы: процесс формирования учебной самоорганизации студентов в системе среднего профессионального образования.

Предмет курсовой работы: методика применения инструмента для составления расписаний, учебных планов и напоминаний как средство повышения учебной самоорганизации студентов ПОО.

В ходе работы использованы методы анализа психолого-педагогической литературы, обобщения педагогического опыта, анкетирование студентов, наблюдение за учебной деятельностью, анализ результатов освоения обучающимися навыков планирования времени. Комплексное использование данных методов позволило всесторонне рассмотреть проблему, определить типичные ошибки студентов при планировании учебного времени и выработать педагогические пути их преодоления.

Практическая значимость курсовой работы заключается в возможности использовать разработанные материалы в педагогической практике колледжей. Результаты исследования будут полезны преподавателям, классным руководителям, методистам, кураторам учебных групп. Предложенная методика позволяет сделать процесс формирования навыков самоорганизации наглядным и практико-ориентированным, способствует адаптации первокурсников к условиям обучения в СПО, повышает общую учебную дисциплину обучающихся.


%==================== ГЛАВА 1 ====================
\newpage
\section{Теоретические основы повышения учебной самоорганизации студентов СПО средствами цифровых инструментов}

\subsection{Сущность и структура понятия учебной самоорганизации студентов}

Современное среднее профессиональное образование ориентировано не только на передачу профессиональных компетенций, но и на формирование самостоятельной, ответственной личности выпускника. В условиях активной цифровизации образовательного процесса от студента требуется умение самостоятельно управлять собственным учебным временем, расставлять приоритеты, контролировать сроки выполнения заданий. Учебная самоорганизация становится важным качеством личности, без которого сложно добиться высоких результатов в обучении и дальнейшей профессиональной деятельности.

В педагогической и психологической литературе учебная самоорганизация рассматривается как интегральное качество личности, которое предполагает сознательное целеполагание, планирование, регулирование и контроль собственной учебной деятельности. Самоорганизация не возникает спонтанно; она формируется постепенно в процессе обучения, под влиянием внешних педагогических воздействий и внутренних усилий самого обучающегося.

Учебная самоорганизация выполняет ряд важных функций. Целеполагающая функция связана с осознанием студентом учебных целей, определением краткосрочных и долгосрочных задач обучения. Планирующая функция предполагает распределение временных ресурсов для достижения поставленных целей. Регулятивная функция отвечает за корректировку действий при возникновении трудностей. Контролирующая функция позволяет студенту самостоятельно оценивать результаты своей работы, выявлять недостатки. Мотивационная функция поддерживает внутреннее стремление обучающегося к достижению учебных результатов.

Структура учебной самоорганизации традиционно описывается через три взаимосвязанных компонента:
\begin{enumerate}
    \item \textbf{Когнитивный компонент.} Включает знания о способах рационального планирования времени, понимание собственных учебных возможностей, знание требований образовательной программы, умение анализировать свои ошибки в организации деятельности.
    \item \textbf{Эмоционально-волевой компонент.} Содержит учебную мотивацию, ответственность, дисциплинированность, способность к самоконтролю, умение преодолевать лень и трудности, устойчивость к отвлекающим факторам.
    \item \textbf{Поведенческий (деятельностный) компонент.} Проявляется в реальных практических действиях: составлении планов и расписаний, соблюдении сроков выполнения заданий, самостоятельной подготовке к занятиям, корректировке собственных планов при изменении условий.
\end{enumerate}

Некоторые исследователи дополнительно выделяют мотивационный компонент, подчеркивая: даже при наличии знаний о правильном планировании студент не будет организовывать свою деятельность без внутреннего желания добиться учебных целей.

Для студентов СПО учебная самоорганизация имеет особое значение. Обучение в колледже сочетает теоретические занятия, лабораторные работы, учебную и производственную практику. У обучающихся появляется большое количество разнообразных заданий с разными сроками сдачи. Отсутствие навыков планирования приводит к перегрузке, прокрастинации, пропускам заданий. Умение работать с расписанием, составлять индивидуальные учебные планы и следовать им готовит студентов к будущей профессиональной деятельности, где от специалиста требуется самостоятельная организация рабочего процесса.

Цифровая среда открывает новые возможности для поддержки самоорганизации. Сервисы для составления расписаний, учебных планов и напоминаний берут на себя часть функций внешнего контроля, помогая студенту фиксировать задачи, отслеживать сроки. Но важно понимать: инструмент сам по себе не формирует самоорганизацию, он является лишь вспомогательным средством. Эффект достигается только тогда, когда обучающийся осваивает принципы планирования и осознанно использует программные возможности.


\subsection{Особенности учебной деятельности и проблемы планирования у студентов СПО}

Период обучения в системе среднего профессионального образования — это важный этап становления личности. Студенты СПО преимущественно находятся в возрастном диапазоне от 15 до 20 лет. Для этого возраста характерна перестройка всей системы учебной деятельности: снижается степень внешнего контроля со стороны педагогов, возрастают требования к самостоятельности обучающегося.

Образовательный процесс в колледже имеет ярко выраженную практическую направленность. Наряду с аудиторными занятиями студенты проходят учебные и производственные практики, выполняют лабораторные, курсовые, индивидуальные проекты. Учебная нагрузка является неоднородной: в отдельные периоды количество заданий резко возрастает (сессия, отчетность по практикам). Это предъявляет повышенные требования к умению распределять время.

В отличие от школы, где расписание и домашние задания контролируются учителями, в СПО значительная часть работы ложится на плечи самого студента. Первокурсники часто не готовы к такой степени самостоятельности, поэтому у них возникают типичные проблемы самоорганизации.

К наиболее распространенным трудностям планирования учебной деятельности студентов СПО относятся:
\begin{enumerate}
    \item Неспособность расставлять приоритеты среди большого количества учебных задач.
    \item Отсутствие навыков составления индивидуального учебного плана поверх основного расписания занятий.
    \item Забывание сроков сдачи отчетов, лабораторных работ, подготовки к зачетам.
    \item Прокрастинация, перенос сложных заданий на последний момент.
    \item Неумение корректировать планы при изменении расписания, переносе занятий, начале практики.
    \item Перегрузка, вызванная нерациональным распределением времени между учёбой, подработкой и досугом.
\end{enumerate}

Атмосфера образовательного учреждения оказывает большое влияние на формирование самоорганизации. Если педагоги уделяют внимание не только передаче профессиональных знаний, но и развитию навыков самостоятельной работы, студенты быстрее осваивают приемы планирования. При авторитарном стиле взаимодействия вся ответственность перекладывается на преподавателя, а у обучающихся не формируется внутренний самоконтроль.

Семейные условия также влияют на уровень самоорганизации. Если в семье не прививались навыки планирования и дисциплины, образовательная организация должна компенсировать этот пробел за счет специально организованной педагогической работы.

Важная роль принадлежит личному примеру преподавателя. Если педагог демонстрирует четкость в соблюдении сроков, использует инструменты планирования, обучающиеся перенимают полезные модели поведения.

Формирование учебной самоорганизации должно охватывать разные направления образовательной работы:
\begin{enumerate}
    \item Учебная деятельность: включение заданий на самостоятельное планирование в содержание дисциплин.
    \item Воспитательная и кураторская работа: групповые занятия-тренинги по управлению временем.
    \item Внеучебная деятельность: опыт организации собственных проектов, участия в студенческих мероприятиях.
\end{enumerate}

Коллективные формы работы и педагогическая рефлексия помогают студентам осознать собственные ошибки в планировании, найти способы их исправления.


\subsection{Педагогические условия применения цифровых инструментов планирования в профессиональных образовательных организациях}

Формирование учебной самоорганизации с использованием цифровых инструментов не происходит само по себе при установке программы на телефон или компьютер. Необходимо создавать комплекс педагогических условий, при которых сервисы для расписаний, планов и напоминаний будут использоваться осознанно, а не формально.

Под педагогическими условиями понимают совокупность организационно-педагогических, психологических и содержательно-методических факторов, обеспечивающих положительный результат педагогического воздействия.

К \textbf{организационно‑педагогическим условиям} относятся:
\begin{itemize}
    \item включение работы с инструментами планирования в планы воспитательной работы кураторов групп;
    \item доступность цифровых сервисов для всех студентов (возможность использования как компьютера, так и мобильного телефона);
    \item системность работы, а не разовые эпизодические занятия. Навыки планирования требуют постоянного повторения и отработки.
\end{itemize}

\textbf{Психологические условия:}
\begin{itemize}
    \item создание благоприятной психологической атмосферы, отсутствие осуждения за ошибки при освоении новых инструментов;
    \item формирование положительной учебной мотивации, демонстрация практической пользы планирования для достижения личных учебных целей;
    \item учет индивидуальных особенностей студентов: разный уровень владения цифровыми технологиями, разная степень сформированности волевых качеств.
\end{itemize}

\textbf{Содержательно‑методические условия:}
\begin{itemize}
    \item использование активных методов обучения: практические упражнения, разбор реальных учебных ситуаций, анализ ошибок планирования;
    \item обучение не только нажатию кнопок в программе, но и теоретическим основам тайм‑менеджмента для студентов;
    \item обязательная рефлексия по итогам использования инструмента: что получилось, какие трудности возникли, что нужно изменить в планах.
\end{itemize}

В работе со студентами СПО важно придерживаться ключевых педагогических принципов.

\textbf{Принцип субъектности}: студент выступает не пассивным пользователем программы, а активным создателем собственных расписаний и учебных планов. Он сам определяет свои задачи, сроки, расставляет приоритеты.

\textbf{Принцип деятельностного подхода}: навыки вырабатываются только в практической деятельности. Недостаточно просто рассказать про возможности инструмента, необходимо, чтобы обучающиеся самостоятельно составили личное расписание, заполнили учебный план, настроили напоминания для своих реальных заданий.

\textbf{Принцип индивидуализации}: студенты имеют разный стартовый уровень. Кто‑то хорошо владеет цифровыми сервисами, а для кого‑то даже простой планировщик является новинкой. Задания должны иметь разную степень сложности.

\textbf{Принцип связи теории с практикой}: все упражнения опираются на реальные задачи конкретной учебной группы: сессия, лабораторные работы, предстоящая производственная практика.

Внеучебная и кураторская работа является важной площадкой для освоения инструментов планирования. Куратор может проводить тематические классные часы, где обучающиеся под руководством педагога осваивают планировщики, делятся опытом их использования.

Среди эффективных педагогических технологий, сочетающихся с применением цифрового инструмента, выделяют:
\begin{itemize}
    \item практикумы‑тренинги по управлению учебным временем;
    \item метод кейсов: анализ ситуаций, когда студент не успел выполнить задание из‑за плохого планирования;
    \item рефлексивные практики, анкетирование для самооценки навыков самоорганизации;
    \item групповое обсуждение опыта использования планировщиков.
\end{itemize}

Эффективность зависит от позиции педагога. Он не должен выступать в роли строгого контролера, который проверяет каждый план студента. Его задача — это наставник, консультант, который показывает возможности инструмента, помогает разобраться с трудностями, подталкивает обучающихся к самостоятельному анализу собственной деятельности.

Важна системность: однократное занятие дает временный эффект. Необходимо периодически возвращаться к теме планирования в течение всего учебного года, усложняя задачи. Сочетание групповых и индивидуальных консультаций позволяет работать со всей группой и отдельно со студентами, имеющими наибольшие трудности с самоорганизацией.


\subsection{Функции инструментов составления расписаний, учебных планов и напоминаний в образовательном процессе}

Инструмент для составления расписаний, учебных планов и напоминаний выполняет несколько взаимосвязанных функций в учебном процессе ПОО. Он не заменяет педагога, но дополняет его работу, выступая вспомогательным средством для становления учебной самоорганизации обучающихся.
\begin{enumerate}
    \item \textbf{Планирующая функция.} Сервис позволяет зафиксировать основное аудиторное расписание, а поверх него формировать индивидуальный учебный план: время для самостоятельной подготовки к занятиям, выполнения домашних и лабораторных заданий, подготовки к зачетам и экзаменам. Студент видит общую картину своей учебной нагрузки, может распределить нагрузку равномерно.
    \item \textbf{Информационно‑напоминающая функция.} Система напоминаний оповещает об истечении сроков сдачи работ, о начале практик, консультаций, зачетов. Это снижает влияние человеческого фактора: забывания важных учебных событий.
    \item \textbf{Регулятивная функция.} При изменении расписания, переносе занятий студент может быстро скорректировать собственный учебный план. Инструмент наглядно показывает перегрузку по дням, помогает перераспределить задачи, чтобы избежать авралов.
    \item \textbf{Контрольно‑рефлексивная функция.} Многие планировщики позволяют отмечать выполненные задачи. Это дает студенту возможность наглядно видеть результаты своей работы, анализировать, какие задания выполняются с опозданием, почему возникли сбои в планах.
    \item \textbf{Мотивационная функция.} Визуализация прогресса выполнения учебного плана повышает удовлетворенность от учебы, укрепляет внутреннюю мотивацию к обучению.
\end{enumerate}

При работе с данным инструментом важно различать два уровня: технический (освоение кнопок, меню, настроек сервиса) и педагогический (освоение принципов грамотного планирования времени). Ошибка многих практик — обучение только техническим операциям, без разбора основ тайм‑менеджмента для студентов. В таком случае обучающиеся заполняют расписание формально, но не используют его в реальной жизни.

Также нужно учитывать риски. Цифровой инструмент может стать источником отвлечения (уведомления от других приложений на телефоне). Педагог должен обсуждать с группой эти опасности, учить обучающихся концентрироваться на задачах.

Таким образом, инструмент для составления расписаний, учебных планов и напоминаний является эффективным вспомогательным педагогическим средством. Его потенциал раскрывается только при грамотной методической организации работы со студентами СПО.


%==================== ГЛАВА 2 ====================
\newpage
\section{Методическая разработка применения инструмента для составления расписаний, учебных планов и напоминаний в работе со студентами СПО}

\subsection{Название, направленность, цели и задачи, форма, методы и средства занятия с использованием инструмента}
\begin{enumerate}
    \item Тема занятия: «Инструмент для составления расписаний, учебных планов и напоминаний — помощник в организации моей учебы».
    \item Направление воспитания: трудовое, интеллектуальное, формирование культуры учебной деятельности.
    \item Форма организации: практическое занятие‑тренинг.
    \item Цель: повышение уровня учебной самоорганизации студентов СПО через освоение практических навыков работы с инструментом для составления расписаний, индивидуальных учебных планов и настройки системы напоминаний.
\end{enumerate}

Задачи мероприятия:
\begin{enumerate}
    \item Сформировать представление обучающихся о значении грамотного планирования учебного времени для успешности обучения в колледже.
    \item Научить студентов практически работать с инструментом: загружать базовое расписание занятий, составлять личный учебный план, расставлять приоритеты задач, настраивать напоминания о сроках сдачи учебных работ.
    \item Отработать умение анализировать собственные ошибки планирования и корректировать индивидуальные учебные планы при изменении учебных условий.
    \item Содействовать формированию ответственности и внутреннего самоконтроля как компонентов учебной самоорганизации будущих специалистов.
\end{enumerate}

Методы воспитания:
\begin{enumerate}
    \item Методы формирования сознания: объяснение, беседа, разбор учебных кейсов.
    \item Методы организации деятельности и формирования поведения: практические упражнения, работа с цифровым инструментом, воспитывающие ситуации.
    \item Методы стимулирования и мотивации: поощрение, демонстрация положительных примеров.
    \item Методы контроля и самоконтроля: самоанализ, анкетирование, рефлексия результатов работы.
\end{enumerate}


\subsection{План практического занятия с указанием этапов и тайминга}

Общая продолжительность занятия: 60 минут.

Оборудование: проектор, экран, мобильные устройства или ноутбуки студентов, презентация‑сопровождение, бланки для анкеты самооценки навыков планирования, пример‑шаблон индивидуального учебного плана.

\textbf{Ход занятия}

\paragraph{1. Организационный этап (10 минут).}
Ведущий (куратор/преподаватель): «Добрый день! Многие из нас сталкивались с ситуацией: мы забываем про срок сдачи лабораторной работы, в последнюю ночь готовимся к зачёту, не успеваем сделать все запланированные дела. Сегодня мы познакомимся с инструментом, который поможет нам упорядочить нашу учебную жизнь. Это инструмент для составления расписаний, учебных планов и напоминаний. Сегодня мы не просто посмотрим на него со стороны — каждый из вас попробует создать собственный реальный учебный план для нашей группы».

Организационный момент: проверка готовности устройств, краткий входной опрос‑беседа. Ведущий задает вопросы группе:
\begin{itemize}
    \item Кто из вас иногда забывает сроки сдачи заданий?
    \item Кто пробовал составлять себе планы учебы? Какие трудности возникали?
\end{itemize}

\paragraph{2. Вводный теоретический этап (10 минут).}
Ведущий рассказывает о сущности учебной самоорганизации, типичных проблемах студентов‑первокурсников, раскрывает основные возможности инструмента: загрузка расписания, создание индивидуальных задач, расстановка приоритетов, напоминания, корректировка планов.

Ведущий: «Наш инструмент не решит все проблемы за нас. Он — только наш помощник. Эффект будет только тогда, когда мы будем осознанно работать с ним, регулярно обновлять свои планы».

Демонстрация на экране основных блоков интерфейса инструмента.

\paragraph{3. Основной практический этап (30 минут).}

Этап 1. Заполнение базового расписания (7 минут).

Ведущий включает таймер: «У вас 7 минут. Откройте инструмент, внесите актуальное расписание наших аудиторных занятий на текущую неделю. Внимательно проверяйте дни и время пар». Педагог ходит по аудитории, помогает студентам, у которых возникают технические трудности, задает наводящие вопросы, но не выполняет задание за обучающихся.

Этап 2. Формирование индивидуального учебного плана (9 минут).

Ведущий: «Теперь поверх расписания пар добавьте ваши личные учебные задачи: подготовка к лабораторным работам, чтение параграфов, подготовка к предстоящему зачету. Расставьте приоритеты: что нужно сделать в первую очередь, что может подождать».

Этап 3. Настройка напоминаний и отработка корректировки плана (12 минут).

Ведущий: «Настройте напоминания о сроках сдачи ваших ближайших заданий. Теперь представим учебную ситуацию‑кейс: завтра одна пара переносится на другой день, добавляется внеплановая консультация. Попробуйте быстро скорректировать ваш учебный план под новые условия». Организуется короткое обсуждение: какие сложности возникли при изменении плана.

Этап 4. Обмен опытом, короткие выступления (2 минуты).

Несколько желающих студентов рассказывают, какие задачи они внесли в свой план, какие напоминания настроили.

\paragraph{4. Рефлексивно‑оценочный этап (10 минут).}
Ведущий: «Практическая часть завершена. Давайте подведем итоги нашей работы. Ответим на несколько вопросов, подумаем над своим опытом».

Вопросы для коллективной рефлексии:
\begin{enumerate}
    \item Что нового вы узнали сегодня про планирование учебного времени?
    \item Какая операция в работе с инструментом оказалась для вас самой сложной?
    \item Как вы планируете использовать полученные навыки в реальной учебе?
    \item Какие опасности при работе с цифровыми планировщиками вы видите?
\end{enumerate}

Студенты заполняют короткий бланк самооценки.

Заключительное слово ведущего: «Сегодня мы освоили важный вспомогательный инструмент. Но умение планировать — это навык, который не появляется после одного занятия. Старайтесь регулярно открывать ваш план, обновлять задачи, следить за напоминаниями. Грамотная самоорганизация помогает не только хорошо учиться в колледже, но и станет очень важной для вашей будущей профессиональной деятельности. Спасибо за работу на занятии».


\subsection{Оценка эффективности применения инструмента для студентов СПО}

Разработанное практическое занятие‑тренинг с применением инструмента для составления расписаний, учебных планов и напоминаний является целостной педагогической методикой, эффективность которой оценивается по ряду критериев, учитывающих специфику среднего профессионального образования.

Главный критерий результативности — продвижение в формировании компонентов учебной самоорганизации обучающихся. В отличие от обычной лекции, где студенты только слушают теорию, данное занятие предполагает полное погружение в практическую деятельность. Обучающиеся работают не с абстрактными примерами, а с реальным расписанием своей учебной группы, реальными предстоящими заданиями. Это повышает личную значимость выполняемых упражнений.

Оценка по поставленным задачам показывает комплексное влияние методики. Образовательная задача по освоению инструмента реализуется на практических этапах: заполнение расписания, создание индивидуальных планов, настройка напоминаний, корректировка планов по условию кейса. Студенты получают не только теоретические сведения, но и отрабатывают конкретные операционные навыки.

Развивающие задачи направлены на формирование умений расставлять приоритеты, анализировать учебную нагрузку, находить выход из изменяющихся учебных ситуаций. Разбор кейса с переносом занятия моделирует типичную ситуацию из реальной жизни колледжа, тренирует гибкость мышления.

Воспитательный эффект проявляется в ходе рефлексии. Студенты осознают собственные трудности с планированием, понимают важность личной ответственности за результаты учебы. Педагог выступает в роли консультанта, не осуждая ошибки обучающихся, что создает комфортные условия для самопознания.

Сценарий занятия ориентирован на аудиторию СПО: задания максимально приближены к реальности обучения колледжа, не перегружены сложной теорией. Занятие подходит как для первокурсников, которые адаптируются к новым условиям, так и для старших курсов при подготовке к сессии и производственной практике.

Косвенными индикаторами эффективности выступают высокая вовлеченность обучающихся в практические задания, активность на этапе рефлексии. Важно понимать: стабильный результат не достигается после одного занятия. Эффект проявляется при повторном обращении к теме планирования в течение семестра. В качестве показателей можно использовать результаты повторного анкетирования самооценки навыков самоорганизации, снижение количества просроченных учебных заданий в группе.

Таким образом, предложенная методика имеет высокий практический потенциал для интеграции в кураторскую и воспитательную работу профессиональных образовательных организаций.


\subsection{Прогноз дальнейшего использования и возможности адаптации инструмента для разных категорий студентов профессиональных образовательных организаций}

Инструмент для составления расписаний, учебных планов и напоминаний и разработанная методика его применения могут выступать универсальным средством для работы с разными группами обучающихся СПО. Методика является гибкой, её элементы легко адаптируются под разные курсы, специальности, уровень цифровой подготовки студентов. Ее можно использовать не только как разовое занятие‑тренинг, но и как элемент долгосрочной воспитательной работы.

Перспективы внедрения в образовательный процесс ПОО. Первое преимущество — доступность. Большинство подобных инструментов работают на мобильных телефонах, которые есть практически у всех студентов. Не требуется сложного специального оборудования.

Возможно использование укороченных версий методики. Например, короткие 5–10‑минутные вставки на классных часах куратора: быстрая проверка и корректировка личных учебных планов, напоминание о предстоящих сроках. Такие мини‑упражнения не требуют большой подготовки, поддерживают навыки планирования на протяжении всего учебного года.

Материал может стать основой целого модуля по формированию культуры учебной деятельности. При регулярных занятиях в течение семестра у обучающихся повышается способность к самостоятельному планированию, уменьшается количество просроченных заданий, спокойнее проходит подготовка к сессии.

Возможна адаптация под разные категории студентов:
\begin{enumerate}
    \item Для первокурсников. Основной акцент делается на адаптацию к обучению в СПО, освоение базовых функций инструмента, знакомство с принципами составления индивидуальных планов.
    \item Для студентов старших курсов. Задания усложняются: добавляется планирование производственной практики, подготовки курсовых и дипломных работ, планирование совмещения учебы и подработки.
    \item Для студентов с низким уровнем владения цифровой техникой. Предусматривается больше индивидуальных консультаций, упрощенные шаблоны планов.
    \item Для групп разных профессиональных направлений содержание кейсов подстраивается под специфику будущей профессии.
\end{enumerate}

При систематическом применении методики можно ожидать следующих устойчивых результатов: повышение уровня учебной самоорганизации обучающихся, снижение количества просроченных заданий, улучшение адаптации первокурсников, формирование важной надпрофессиональной компетенции планирования времени, востребованной на рынке труда.


%==================== ЗАКЛЮЧЕНИЕ ====================
\newpage
\section*{Заключение}
\addcontentsline{toc}{section}{Заключение}

В ходе курсовой работы было рассмотрено педагогическое значение инструмента для составления расписаний, учебных планов и напоминаний как средства повышения учебной самоорганизации студентов профессиональных образовательных организаций. Учебная самоорганизация выступает важнейшей интегральной характеристикой личности будущего специалиста среднего звена, она включает когнитивный, эмоционально‑волевой и поведенческий компоненты и не формируется спонтанно.

Теоретический анализ показал, что у студентов СПО существует ряд типичных трудностей планирования учебного времени: неумение расставлять приоритеты, забывание сроков сдачи работ, прокрастинация, сложности с корректировкой планов при изменении учебных условий. Цифровые инструменты планирования могут существенно помочь в решении этих проблем, но только при создании комплекса педагогических условий. Важно различать техническое освоение сервиса и педагогическую работу по формированию навыков осознанного планирования времени.

Итогом теоретического исследования стала разработанная методика практического занятия‑тренинга. Сценарий рассчитан на стандартные 60 минут академического времени, подходит для проведения куратором группы или преподавателем. Занятие сочетает короткий теоретический блок, практическую работу с реальными учебными задачами группы, разбор педагогического кейса и обязательный рефлексивный этап.

Разработанная методика обладает гибкостью: её можно адаптировать для первокурсников и старших курсов, использовать в полном виде или в формате коротких упражнений на классных часах. Практическая ценность заключается в том, что все упражнения опираются на реальную учебную ситуацию конкретной студенческой группы. Рефлексия переводит полученный практический опыт в осознанные личные выводы студентов.

Внедрение подобных методик в повседневную работу колледжей способствует повышению уровня учебной самоорганизации обучающихся, улучшению их академических результатов, формированию надпрофессиональных компетенций, которые понадобятся выпускникам в дальнейшей трудовой деятельности. Формирование учебной самоорганизации возможно только при сочетании теоретических знаний, практической деятельности с цифровыми инструментами и целенаправленной педагогической поддержки со стороны преподавателей и кураторов.

Инструмент для составления расписаний, учебных планов и напоминаний не решает все педагогические задачи сам по себе, но выступает эффективным вспомогательным средством, связывающим учебный процесс и личностное развитие будущих специалистов СПО.


%==================== СПИСОК ЛИТЕРАТУРЫ ====================
\newpage
\section*{Список литературы}
\addcontentsline{toc}{section}{Список литературы}

\begin{enumerate}
    \item Андреева, Г. М. Психология общения : учебное пособие для вузов / Г. М. Андреева. — Москва : Аспект Пресс, 2012. — 368 с.
    \item Бодалев, А. А. Психология общения : избранные психологические труды / А. А. Бодалев. — 4‑е изд., перераб. и доп. — Москва : МПСИ ; Воронеж : МОДЭК, 2008. — 318 с.
    \item Коджаспирова, Г. М. Педагогика в схемах и таблицах : учебное пособие / Г. М. Коджаспирова. — Москва : Проспект, 2016. — 248 с.
    \item Леонтьев, А. А. Педагогическое общение / А. А. Леонтьев. — 2‑е изд., перераб. и доп. — Москва : НИИ школьных технологий, 2009. — 368 с.
    \item Петровская, Л. А. Компетентность в общении: социально‑психологический тренинг / Л. А. Петровская. — Москва : Изд‑во Моск. ун‑та, 1989. — 216 с.
    \item Захарова, И. Г. Цифровые инструменты в организации самостоятельной работы студентов СПО / И. Г. Захарова // Среднее профессиональное образование. — 2021. — № 6.
\end{enumerate}


%==================== ПРИЛОЖЕНИЕ ====================
\newpage
\section*{Приложение}
\addcontentsline{toc}{section}{Приложение}

\textbf{Приложение 1}\\
Шаблон индивидуального учебного плана (пример для заполнения в инструменте)

\begin{tabular}{|l|l|l|l|l|l|}
\hline
Период & Аудиторные занятия & Самостоятельные учебные задачи & Срок сдачи & Напоминание & Приоритет \\
\hline
Понедельник & Расписание пар & Подготовка к лабораторной работе №2 & Среда 14:00 & За день до сдачи & Высокий \\
\hline
Вторник & Расписание пар & Прочитать параграф учебника & Пятница & За два дня & Средний \\
\hline
\end{tabular}

\end{document}
